% Fichier généré automatiquement par tools/generate_corrections.py.
% Source : STAT/STAT-03-Demarche/04_RR/04_RR.tex
% Statut : complet (4/4 question(s) renseignée(s), 0 reprise(s)).
\begin{corrigebox}[Corrigé extrait de la banque]
\CorrigeQuestion{1}
\begin{marginfigure}
\includegraphics[width=\linewidth]{04_RR_01_cor}
\end{marginfigure}
\CorrigeQuestion{2}
\begin{itemize}
\item Pivot entre 0 et 1 : $\torseurstat{T}{0}{1}=\torseurl{X_{01}\vi{0}+Y_{01}\vj{0}+Z_{01}\vk{0}}{M_{01}\vj{0}+N_{01}\vk{0}}{A,\rep{0}}$.
\item Pivot entre 1 et 2 : $\torseurstat{T}{1}{2}=\torseurl{X_{12}\vi{1}+Y_{12}\vj{1}+Z_{12}\vk{1}}{M_{12}\vj{1}+N_{12}\vk{1}}{B,\rep{0}}$.
\item Pesanteur sur 1: $\torseurstat{T}{\text{pes}}{1}=\torseurl{-m_1g\vj{0}}{\vect{0}}{G_1,\rep{0}}$.
\item Pesanteur sur 2: $\torseurstat{T}{\text{pes}}{2}=\torseurl{-m_2g\vj{0}}{\vect{0}}{G_2,\rep{0}}$.
\item Moteur entre 0 et 1 : $\torseurstat{T}{0_{m1}}{1}=\torseurl{\vect{0}}{C_1\vk{0}}{A,\rep{0}}$.
\item Moteur entre 1 et 2 : $\torseurstat{T}{1_{m2}}{2}=\torseurl{\vect{0}}{C_2\vk{0}}{B,\rep{0}}$.
\end{itemize}
\CorrigeQuestion{3}
\begin{itemize}
\item Pivot entre 0 et 1 : $\torseurstat{T}{0}{1}=\torseurl{X_{01}\vi{0}+Y_{01}\vj{0}}{\vect{0}}{A,\rep{0}}$.
\item Pivot entre 1 et 2 : $\torseurstat{T}{1}{2}=\torseurl{X_{12}\vi{1}+Y_{12}\vj{1}}{\vect{0}}{B,\rep{0}}$.
\item Pesanteur sur 1: $\torseurstat{T}{\text{pes}}{1}=\torseurl{-m_1g\vj{0}}{\vect{0}}{G_1,\rep{0}}$.
\item Pesanteur sur 2: $\torseurstat{T}{\text{pes}}{2}=\torseurl{-m_2g\vj{0}}{\vect{0}}{G_2,\rep{0}}$.
\item Moteur entre 0 et 1 : $\torseurstat{T}{0_{m1}}{1}=\torseurl{\vect{0}}{C_1\vk{0}}{A,\rep{0}}$.
\item Moteur entre 1 et 2 : $\torseurstat{T}{1_{m2}}{2}=\torseurl{\vect{0}}{C_2\vk{0}}{B,\rep{0}}$.
\end{itemize}
\CorrigeQuestion{4}
C'est une chaîne ouverte. On isole l'extrémité et on applique le théorème correspondant aux mobilités : 
\begin{itemize}
\item on isole \textbf{2} et on réalise le théorème du moment statique en $A$ en projection sur $\vk{0}$;
\item on isole \textbf{1+2} et on réalise le théorème du moment statique en $B$ en projection sur $\vk{0}$.
\end{itemize}
\end{corrigebox}
